\documentclass[11pt]{article}
\usepackage[margin=1in]{geometry}
\usepackage{amsmath, amssymb}
\usepackage{enumitem}
\setlist[itemize]{itemsep=2pt, topsep=3pt}

\title{Compressed Computation can be Computation in Superposition}
\author{Francisco Ferreira da Silva}
\date{\today}

\begin{document}
\maketitle

\section*{Thesis}

Computation in superposition (CiS) is a central hypothesis in mechanistic interpretability.
Braun et al. introduce a compressed computation (CC) toy model that appeared to implement CiS.
However, this was shown by Bhagat et al.\ to most likely not be a toy model for CiS.
To the best of our knowledge, no suitable toy model of CiS has been introduced.

We provide one.
We consider a single-layer network with fewer ReLU neurons than features and train it under $L^4$ loss to learn the ReLU of its inputs.
This model is identical to that of Bhagat et al., with the exception of the loss we choose.
By penalizing outliers more harshly, this loss encourages CiS.
We argue our model exhibits CiS, and we reverse engineer its learned mechanism to support this.
We claim it implements \emph{combinatorial ReLU coding}: the $n$ bottleneck neurons partition input space into $2^n - 1$ nontrivial sign regions, each acting as a codeword; the trained model assigns features to codewords, and within each region an effective low-rank linear map reads the assigned features out. Two distinct bottlenecks govern performance (codeword identification and rank-$n$ readout); we characterise both and separate them experimentally.

\section{Introduction}

\paragraph{What the paper claims.}
\begin{enumerate}
\item After the Bhagat et al.\ critique, the field does not have a clean toy model of CiS.
\item Training this architecture under $L^4$ loss yields a model that genuinely does CiS.
\item The learned mechanism is combinatorial ReLU coding, which we reverse engineer in the smallest nontrivial case and predict quantitatively at larger sizes.
\item The apparent-CiS vs.\ genuine-CiS distinction has a clean mechanistic handle: whether the model exploits the $2^n - 1$ sign regions of the ReLU bottleneck.
\end{enumerate}

\paragraph{Context.}
\begin{itemize}
\item Anthropic's Toy Models of Superposition established that superposed \emph{representations} are robust in small models.
\item Whether \emph{computation} composes under superposition is the natural follow-up and is under active discussion (Braun et al.; Linsefors/Bushnaq on multi-layer noise propagation; Bhagat et al.'s critique).
\item Our contribution sits directly in that conversation: a post-Bhagat replacement toy model, plus a mechanistic account.
\end{itemize}

\section{Why $L^4$ loss}

Under $L^2$ loss, the naive solution (implement the first $n$ features exactly and ignore the remaining $F - n$) is what optimisation finds.
The reason is that $L^2$ penalises average squared error, and with sparse inputs the average is dominated by the many features that are rarely active.
Getting a few features exactly right and ignoring the rest scores well on this objective, and no gradient pressure exists to do better.

$L^4$ loss changes this.
By raising the error exponent, the loss disproportionately penalises the worst-served features.
Features that are completely ignored now cost more than features that are served imperfectly.
This creates gradient pressure to assign some representational capacity to every feature, which is the prerequisite for CiS.

We support this picture with two pieces of evidence.
First, loss-per-feature plots for $L^2$ vs.\ $L^4$: under $L^2$ the distribution is bimodal (matching naive), under $L^4$ it is approximately uniform.
Second, an exponent sweep showing a crossover: as the exponent increases past $2$, the trained model's total loss crosses below the naive baseline and the per-feature loss distribution flattens.

\paragraph{Controls we commit to running.}
\begin{itemize}
\item \emph{$L^4$-on-naive-class control.} Constrain the model to the naive solution class and fit under $L^4$.
If $L^4$ alone, without CiS, already matches the trained model's loss, then the exponent change is not actually eliciting new behaviour.
We predict it does not.
\item \emph{Sparsity sweep.} Vary $p$ at fixed $F, n$.
Identify where the trained solution departs from naive and where it collapses back.
\item \emph{Task-structure sweep.} Repeat with target functions other than $\mathrm{ReLU}(x)$: e.g.\ $x^2$, $|x|$, thresholded ReLU.
If the same mechanism appears, the toy model is not specific to the near-identity task.
\end{itemize}

\section{Evidence that the model does CiS}

\begin{itemize}
\item \textbf{Total loss strictly beats the naive baseline.}
The architecture has no component capable of the Braun-style mixing that Bhagat et al.\ showed was responsible for the appearance of CiS in prior work, so the improvement is attributable to the bottleneck itself.
\item \textbf{Loss is approximately uniform across features.}
The naive solution has zero loss on the first $n$ features and unit loss on the rest; the CiS solution spreads error roughly evenly.
\item \textbf{Frame geometry.} $W_{\mathrm{in}}$ columns form an approximate tight frame.
Baselines: random $W_{\mathrm{in}}$ and $L^2$-trained $W_{\mathrm{in}}$.
\item \textbf{Interference signature.} Per-sample loss grows with the number of simultaneously active features.
\item \textbf{Scaling with $n$ at fixed expansion ratio.}
One neuron cannot do CiS; two can; the gap to naive grows with $n$.
\end{itemize}
These establish the phenomenon.
Section~4 explains it.

\section{The mechanism: combinatorial ReLU coding}

\paragraph{Codewords from ReLU sign regions.}
The $n$ preactivation hyperplanes partition $\mathbb{R}^n$ into up to $2^n$ orthants.
The all-negative orthant produces zero output; the remaining $2^n - 1$ regions each correspond to a distinct subset of active neurons.
We call this subset the \emph{codeword} of the input.

\paragraph{Feature-to-codeword assignment.}
Each feature $j$ has an input direction $c_j = W_{\mathrm{in}}[:,j]$.
When $x = v e_j$ with $v > 0$, the codeword fired is the sign pattern of $c_j$.
In the trained model, the columns of $W_{\mathrm{in}}$ cluster onto the $2^n - 1$ codewords, partitioning the $F$ features into groups.

\paragraph{Per-region effective map.}
Within the region where codeword $s$ fires, the network acts as the linear map $x \mapsto W_{\mathrm{out}} \, \mathrm{diag}(s) \, W_{\mathrm{in}} x$.
Different regions give different effective maps, and each region is optimised primarily for the features assigned to its codeword.

\paragraph{Reverse engineering the $n=2$, $F=20$ case.}
Symmetries (B$\leftrightarrow$C mirror, A neuron-swap, ReLU homogeneity) reduce the 12 free weights to a 5-continuous + 1-discrete parameter ansatz.
Numerically the ansatz matches the trained model's $L^4$ loss within MC noise at 2M samples.
This is a full mechanistic account: closed-form weights, closed-form loss, numerical match.

\paragraph{The two bottlenecks.}
\begin{itemize}
\item \textbf{Identification} (which feature is active): capped by the number of codewords, $2^n - 1$.
When $F > 2^n - 1$, features must share codewords and single-feature inputs are fundamentally unidentifiable.
\item \textbf{Readout} (recovering the value given identification): $W_{\mathrm{out}}$ has rank $\leq n$, so the set of reachable outputs is at most $n$-dimensional.
Perfect reconstruction of $F$ orthogonal targets requires $F \leq n$; otherwise there is a Frobenius lower bound of $F - n$ on the single-feature response matrix's distance from identity.
\end{itemize}

\paragraph{Separation experiment (MLP decoder).}
Replacing $W_{\mathrm{out}}$ with a nonlinear decoder keeps identification bounded by codewords but lifts the rank constraint.
This isolates the two bottlenecks: at $n$ large enough for unique codewords, single-feature loss should drop to near zero; at $n$ too small, it should plateau at a codeword-ambiguity floor.
Reported side-by-side with the linear-decoder numbers.

\section{Predictions and verification}

The mechanism supplies falsifiable quantitative predictions.
For each we state the metric and what would constitute failure.

\begin{enumerate}
\item \textbf{Feature partition.} For every $n$, features cluster onto at most $2^n - 1$ codewords.
Metric: fraction of features whose codeword is shared with others, and expected partition sizes.
Failure mode: features using more than $2^n - 1$ distinct codewords, or failing to concentrate onto any discrete set.
\item \textbf{Phase transition at $F = 2^n - 1$.}
Fixing $F$ and increasing $n$, single-feature $L^4$ loss drops sharply at the $n$ where codewords first become unique.
Example: $F = 20$, transition at $n = 5$.
\item \textbf{Rank-$n$ readout bound.}
For $F > n$, the distance of the scaled single-feature response matrix from $I_F$ approaches the information-theoretic floor $F - n$.
Stated across $n = 3, 4, 5$, not just $n = 5$.
\item \textbf{Structured assignment.}
The trained codeword assignment is structured, not a random injection.
Control: shuffle assigned codewords, retrain the decoder only, compare loss.
\item \textbf{Sparsity regime.}
Codeword structure persists over an explicit range of $p$; we identify the breakdown boundary.
\item \textbf{Seed robustness.}
The partition structure, not the exact per-feature assignment, emerges reliably across seeds.
If the ansatz has multiple basins (the discrete parameter suggests it might), we report frequencies.
\end{enumerate}

\section{Planned experiments (working back from claims)}

This section lists what we will run, organised by which claim each experiment supports.
Not all are in the current run.

\begin{itemize}
\item \textbf{Against $L^4$-as-magic:} $L^4$-on-naive-class control; exponent sweep over $\{1.5, 2, 3, 4, 6, 8\}$.
\item \textbf{Against ``ReLU(x) is special":} alternative targets ($x^2$, $|x|$, thresholded ReLU, simple 2-feature products).
\item \textbf{Against ``only works at our $p$":} sparsity sweep $p \in [0.005, 0.3]$.
\item \textbf{For generality of the codeword story:} larger $F, n$ (e.g.\ $F = 100$, $n \in \{2, \ldots, 10\}$); ansatz extension to $n = 3$ (7 groups).
\item \textbf{For separation of bottlenecks:} MLP decoder at multiple $n$; rank-constrained decoder variants.
\item \textbf{For robustness of the story:} seed sweep with basin frequencies; initialisation sweep.
\item \textbf{Negative controls:} random-codeword-assignment ablation; $L^2$-trained baseline on every plot.
\end{itemize}

\section{Discussion}

\paragraph{Relation to Bhagat et al.}
Their critique showed that prior apparent-CiS behaviour was attributable to a component of the architecture that our model does not contain.
The mechanism we identify here is a property of the ReLU bottleneck's sign-region geometry, which was present in their analysis but not doing the work.

\paragraph{Relation to Linsefors / Bushnaq.}
Their analysis of multi-layer noise propagation is complementary.
The codeword mechanism solves identification at a single layer.
How identification errors compose across layers is the natural next question and is outside the scope of this paper.

\paragraph{Relation to Anthropic's Toy Models of Superposition.}
Those papers establish that \emph{representations} superpose.
We give a minimal model in which \emph{computation} does, and show what limits it.

\paragraph{Limitations.}
Single layer.
Near-identity task (ReLU of input), though the task-structure sweep partially addresses this.
Small scale.
Whether combinatorial ReLU coding is observable in natural networks is open.

\paragraph{What we are not claiming.}
We are not claiming that combinatorial ReLU coding is the unique CiS mechanism, nor that it scales to realistic networks unchanged.
The claim is that it is a clean, fully-characterised mechanism in the smallest setting where genuine CiS occurs, and that it answers the question Bhagat et al.\ left open.

\end{document}
